\section{Amplifier ou réduire un quotient}

\begin{aconnaitre}
\textbf{Si on multiplie ou si on divise} le numérateur et le dénominateur d'un quotient par \textbf{un même nombre non nul} alors on obtient \textbf{quotient égal}.

Pour tous nombres $a$, $b$ et $k$ où $b$ et $k$ sont non nuls :
\[ \frac{a \times k}{b \times k} = \frac{a}{b} \text{ et } \frac{a \div k}{b \div k} = \frac{a}{b}\]
\end{aconnaitre}

\begin{exemple*1}
Simplifie le quotient $\dfrac{42}{-140}$.
\correction

\begin{tabular}{lll}
$\dfrac{42}{-140}=-\dfrac{42}{140}$ & $\longrightarrow$ & On détermine le signe du quotient. \\
$\dfrac{42}{-140}=-\dfrac{3 \times 2 \times 7}{10 \times 7 \times 2}$  & $\longrightarrow$ & On cherche les facteurs communs à 2 et 140. \\
$\dfrac{42}{-140}=-\dfrac{3}{10}$ & $\longrightarrow$ & On simplifie le quotient. \\
\end{tabular}
\end{exemple*1}

\begin{exemple*1}
Détermine le nombre manquant dans l'égalité $\dfrac{-1,2}{6}=\dfrac{...}{18}$.
\correction

\vspace{.5em}
\begin{tabular}{lll}
\hspace{1.5em}\includegraphics[width=.8cm]{figX3} & & \\
$\dfrac{-1,2}{6}=\dfrac{...}{18}$ & $\longrightarrow$ & Pour passer de 6 à 18 \textbf{on multiplie par 3}. \\
\hspace{1.5em}\includegraphics[width=.8cm]{figX3b} & & \\
donc $\dfrac{-1,2}{6}=\dfrac{-3,6}{18}$  & $\longrightarrow$ & Ainsi, pour trouver le nombre manquant, \textbf{on multiplie} \\
& & \textbf{$\mathbf{-1,2}$ par 3}, ce qui donne $-3,6$. \\
\end{tabular}
\end{exemple*1}


\section{Comparer deux fractions}

\begin{aconnaitre}[Fractions dont le numérateur et le dénominateur sont positifs]
\begin{itemize}
\item Si les deux fractions ont le même dénominateur, la plus grande est celle qui a le plus grand numérateur.

Par exemple $\dfrac{5}{7}>\dfrac{2}{7}$ car $5 > 2$.

\item Si les deux fractions ont le même numérateur, la plus grande est celle qui a le plus petit dénominateur.

Par exemple $\dfrac{4}{9}<\dfrac{4}{5}$ car $9 > 5$.

\item Si une fraction est plus grande que 1 (son numérateur est plus grand que son dénominateur), alors elle est plus grande qu'une fraction qui est plus petite que 1 (dont le numérateur est plus petit que le dénominateur).

Par exemple $\dfrac{9}{5}>\dfrac{6}{7}$ car $\dfrac{9}{5}>1$ et $\dfrac{6}{7}<1$. 
\end{itemize}
\end{aconnaitre}



\begin{remarque}
Attention, les règles données ci-dessus ne sont pas vraies si le numérateur ou le dénominateur d'une des fractions est négatif.
\end{remarque}

\section{Multiplication}

\begin{aconnaitre}
Pour \textbf{multiplier des nombres en écriture fractionnaire}, on multiplie les numérateurs entre eux et les dénominateurs entre eux.

Pour tous nombres $a$, $b$, $c$ et $d$ où $b$ et $d$ sont non nuls :
\[ \dfrac{a}{b} \times \dfrac{c}{d} = \dfrac{a\times c}{b \times d} \]
\end{aconnaitre}

\begin{remarque}
Si $b=1$, la formule devient $a \times \dfrac{c}{d} = \dfrac{a\times c}{d}$
\end{remarque}

\begin{exemple*1}
Calcule l'expression $B=-\dfrac{35}{33} \times \dfrac{-39}{-80}$. Donne le résultat sous forme simplifiée.
\correction
\vspace{.5em}
 
\begin{tabular}{lll}
$B = -\dfrac{35 \times 39}{33 \times 80}$ & $\longrightarrow$ & On détermine le signe du résultat. \\
$B = -\dfrac{7 \times \mathbf{5} \times 13 \times \mathbf{3}}{11 \times \mathbf{3} \times 2 \times \mathbf{5} \times 8}$ & $\longrightarrow$ & On cherche des facteurs communs. \\
$B = -\dfrac{7 \times 13}{11 \times 2 \times 8}$ & $\longrightarrow$ & On simplifie. \\
$B = -\dfrac{91}{176}$ & $\longrightarrow$ & On calcule. \\
\end{tabular}
\end{exemple*1}

\section{Division de deux quotients}

\subsection{Inversion d'un nombre non nul}

\begin{definition}
\textbf{Deux nombres sont \MotDefinition{inverses}{}} l'un de l'autre si leur produit est égal à 1.
\end{definition}

\begin{propriete}
\begin{itemize}
    \item Tout nombre $x$ non nul admet un inverse qui est le nombre $\dfrac{1}{x}$.
    \item Tout nombre en écriture fractionnaire $\dfrac{a}{b}$ ($a\neq 0$ et $b \neq 0$) admet un inverse qui est le nombre $\dfrac{b}{a}$.
\end{itemize}
\end{propriete}

\begin{remarque}
\begin{itemize}
    \item Un nombre et son inverse ont toujours le même signe.
    
    En effet, leur produit 1 est positif et seul le produit de deux nombres de même signe est positif.
    \item Zéro est le seul nombre qui n'admet pas d'inverse.
    
    En effet, tout nombre multiplié par 0 donne 0 et ne donnera jamais 1.
\end{itemize}
\end{remarque}

\begin{exemple}
Calcule l'inverse de 3.
\correction
L'inverse de 3 est $\dfrac{1}{3}$.
\end{exemple}

\begin{exemple}
Calcule l'inverse de $\dfrac{-7}{3}$.
\correction
L'inverse de $\dfrac{-7}{3}$ est $\dfrac{1}{\dfrac{-7}{3}}=\dfrac{3}{-7}=\dfrac{-3}{7}$.
\end{exemple}

\subsection{Diviser des quotients}

\begin{aconnaitre}
\textbf{Diviser par un nombre non nul} revient à multiplier par l'inverse de ce nombre.

Pour tous nombres $a$, $b$, $c$ et $d$ où $b$, $c$ et $d$ sont non nuls :
\[ \dfrac{a}{b} \div \dfrac{c}{d} = \dfrac{a}{b} \times \dfrac{d}{c} \text{ ou } \dfrac{\phantom{i}\dfrac{a}{b}\phantom{i}}{\dfrac{c}{d}}=\dfrac{a}{b} \times \dfrac{d}{c} \]
\end{aconnaitre}

\begin{exemple*1}
Calcule $C=\dfrac{-8}{7} \div \dfrac{5}{-3}$.
\correction
\vspace{.5em}
 
\begin{tabular}{lll}
$C=+\left(\dfrac{8}{7} \div \dfrac{5}{3}\right)$ & $\longrightarrow$ & On détermine le signe du résultat. \\
$C=\dfrac{8}{7} \times \dfrac{3}{5}$ & $\longrightarrow$ & On multiplie par l'inverse du deuxième quotient. \\
$C=\dfrac{8\times 3}{7\times 5}$ & $\longrightarrow$ & On multiplie les fractions. \\
$C=\dfrac{24}{35}$ & $\longrightarrow$ & On calcule. \\
\end{tabular}
\end{exemple*1}


\begin{exemple*1}
Calcule $D=\dfrac{-\dfrac{32}{21}}{\dfrac{-48}{-35}}$ et donne le résultat en le simplifiant le plus possible.
\correction
\vspace{.5em}
 
\begin{tabular}{lll}
$D=-\dfrac{\dfrac{32}{21}}{\dfrac{48}{35}}$ & $\longrightarrow$ & On détermine le signe du résultat. \\
$D=-\dfrac{32}{21}\times \dfrac{35}{48}$ & $\longrightarrow$ & On multiplie par l'inverse du deuxième quotient. \\
$D=-\dfrac{\mathbf{8}\times \mathbf{2}\times 2 \times \mathbf{7} \times 5}{\mathbf{7 \times 3 \times 3 \times \mathbf{2} \times \mathbf{8}}}$ & $\longrightarrow$ & On cherche des facteurs communs. \\
$D=-\dfrac{10}{9}$ & $\longrightarrow$ & On calcule sans oublier de simplifier avant ! \\
\end{tabular}
\end{exemple*1}


\begin{exemple*1}
Quelle est la nature du nombre $E$ défini par $E=\dfrac{1+\dfrac{2}{3}}{1-\dfrac{2}{3}}$ ?
\correction
\vspace{.5em}
 
\begin{tabular}{lll}
$E=\dfrac{\dfrac{3}{3}+\dfrac{2}{3}}{\dfrac{3}{3}-\dfrac{2}{3}}=\dfrac{\dfrac{5}{3}}{\dfrac{1}{3}}$ & $\longrightarrow$ & $E$ peut s'écrire aussi $E=\left(1+\dfrac{2}{3}\right) \div \left(1-\dfrac{2}{3}\right)$. \\
 & & On commence donc par calculer les parenthèses.\\
$E=\dfrac{5}{3} \times \dfrac{3}{1}$ & $\longrightarrow$ & On multiplie par l'inverse du deuxième quotient. \\
$E=\dfrac{5\times \mathbf{3}}{\mathbf{3}\times 1}$ & $\longrightarrow$ & On cherche des facteurs communs. \\
$E=5$ donc $E$ est un nombre entier. & &  \\
\end{tabular}
\end{exemple*1}


\section{Signe d'une fraction}

\begin{propriete}
Pour déterminer le signe d'une fraction, on utilise la règle du produit des signes.
\end{propriete}

\begin{aconnaitre}
Si une fraction est négative, on peut l'écrire de trois manières différentes en mettant le signe moins au numérateur, au dénominateur (peu utilisé) ou devant la fraction.
On écrira donc indifféremment  $\dfrac{-2}{3}$ ou $-\dfrac{2}{3}$ et rarement $\dfrac{2}{-3}$.
\end{aconnaitre}

